% GENERATED FILE. Edit canonical lessons, then run npm run generate:books.
% canonical-source-hash: fnv1a64:217b180bc9ba4bdc
% canonical-lessons: GU-C497-kankotri, GU-C497-pithi, GU-C497-jan, GU-C497-varghodo, GU-C497-parnavvu

\chapter{Wedding Plans}
\label{ch:gu-wedding-plans}
\hlchaptermodality{497}

\begin{chapteropening}
\textbf{By the end of this chapter:} \emph{I can talk about getting ready for a wedding.}

Complete the last lesson of chapter 497: I can talk about getting ready for a wedding.
\end{chapteropening}

\section[\texorpdfstring{\gu{કંકોતરી} (kankotrī)}{કંકોતરી (kankotrī)}]{\gu{કંકોતરી} (kankotrī) — a wedding invitation card}

\label{lesson:GU-C497-kankotri}



\begin{quote}
Before the new one: say the Gujarati for to welcome with rites, then the Gujarati for a baby's nappy.
\end{quote}

\subsection*{You'll want to know: \gu{કંકોતરી}}
\textbf{\gu{કંકોતરી}} — \emph{kankotrī} — \textquotedblleft{}a wedding invitation card\textquotedblright{}.

\begin{quote}
\gu{અમે} \gu{બધા} \gu{સંબંધીઓને} \gu{કંકોતરી} \gu{મોકલી}.
\end{quote}

\begin{quote}
\emph{ame badhā sambandhīone kankotrī moklī.}
\end{quote}

\begin{quote}
We sent wedding invitations to all our relatives.
\end{quote}

A feminine noun.

\begin{cousinweb}
\textbf{\gu{કંકોતરી}} (\emph{kankotrī}) goes back to a reconstructed \emph{*kuṅkumapattrikā-}, \textquotedblleft{}a saffron letter\textquotedblright{}, from \emph{kuṅkuma}, \textquotedblleft{}saffron\textquotedblright{}, and \emph{pattrikā}, \textquotedblleft{}letter\textquotedblright{}; CDIAL lists the Gujarati word, an invitation dyed red (CDIAL 3215).
\end{cousinweb}

\begin{grammarlens}[title={What you've built}]
One more word for describing what you lived and felt.
\end{grammarlens}

\subsection*{Your turn}
\begin{itemize}
\raggedright
  \item \emph{Say it:} \emph{kankotrī}
  \item \emph{Say it:} \emph{kankotrī}, once more
  \item \emph{Say it:} \emph{ame badhā sambandhīone kankotrī moklī.} — We sent wedding invitations to all our relatives.
  \item \emph{Recall:} say the Gujarati for to welcome with rites, then the Gujarati for a baby's nappy, then say \emph{kankotrī} again
\end{itemize}

\begin{tcolorbox}[breakable,colback=teal!4,colframe=teal!35!black,title={Before you move on}]
What does \gu{કંકોતરી} mean? (\textquotedblleft{}A wedding invitation card\textquotedblright{}.) Say it once more.
\end{tcolorbox}
\section[\texorpdfstring{\gu{પીઠી} (pīṭhī)}{પીઠી (pīṭhī)}]{\gu{પીઠી} (pīṭhī) — the turmeric rite; turmeric paste put on before a wedding}

\label{lesson:GU-C497-pithi}



\begin{quote}
Before the new one: say the Gujarati for a baby's nappy, then the Gujarati for a wedding invitation card.
\end{quote}

\subsection*{You'll want to know: \gu{પીઠી}}
\textbf{\gu{પીઠી}} — \emph{pīṭhī} — \textquotedblleft{}the turmeric rite; turmeric paste put on before a wedding\textquotedblright{}.

\begin{quote}
\gu{લગ્ન} \gu{પહેલાં} \gu{બહેનપણીઓએ} \gu{મીરાને} \gu{પીઠી} \gu{લગાડી}.
\end{quote}

\begin{quote}
\emph{lagna pahelā̃ bahenpaṇīoe Mīrāne pīṭhī lagāḍī.}
\end{quote}

\begin{quote}
Before the wedding Mira's friends put turmeric paste on her.
\end{quote}

A feminine noun.

\begin{cousinweb}
\textbf{\gu{પીઠી}} (\emph{pīṭhī}) is inherited from Sanskrit \emph{piṣṭa}, \textquotedblleft{}crushed, ground\textquotedblright{}; CDIAL lists the Gujarati word, the yellow powder put on bride and groom (CDIAL 8218).
\end{cousinweb}

\begin{grammarlens}[title={What you've built}]
One more word for describing what you lived and felt.
\end{grammarlens}

\subsection*{Your turn}
\begin{itemize}
\raggedright
  \item \emph{Say it:} \emph{pīṭhī}
  \item \emph{Say it:} \emph{pīṭhī}, once more
  \item \emph{Say it:} \emph{lagna pahelā̃ bahenpaṇīoe Mīrāne pīṭhī lagāḍī.} — Before the wedding Mira's friends put turmeric paste on her.
  \item \emph{Recall:} say the Gujarati for a baby's nappy, then the Gujarati for a wedding invitation card, then say \emph{pīṭhī} again
\end{itemize}

\begin{tcolorbox}[breakable,colback=teal!4,colframe=teal!35!black,title={Before you move on}]
What does \gu{પીઠી} mean? (\textquotedblleft{}The turmeric rite; turmeric paste put on before a wedding\textquotedblright{}.) Say it once more.
\end{tcolorbox}
\section[\texorpdfstring{\gu{જાન} (jān)}{જાન (jān)}]{\gu{જાન} (jān) — the groom's wedding party}

\label{lesson:GU-C497-jan}



\begin{quote}
Before the new one: say the Gujarati for a wedding invitation card, then the Gujarati for the turmeric rite.
\end{quote}

\subsection*{You'll want to know: \gu{જાન}}
\textbf{\gu{જાન}} — \emph{jān} — \textquotedblleft{}the groom's wedding party\textquotedblright{}.

\begin{quote}
\gu{સાંજે} \gu{જાન} \gu{ગામમાં} \gu{આવી}.
\end{quote}

\begin{quote}
\emph{sā̃je jān gāmmā̃ āvī.}
\end{quote}

\begin{quote}
In the evening the groom's wedding party arrived in the village.
\end{quote}

A feminine noun.

\begin{cousinweb}
\textbf{\gu{જાન}} (\emph{jān}) is inherited from Sanskrit \emph{janya}, \textquotedblleft{}belonging to a people\textquotedblright{} and \textquotedblleft{}a friend of the bridegroom\textquotedblright{}; CDIAL lists the Gujarati word, \textquotedblleft{}bridegroom's procession\textquotedblright{} (CDIAL 5116).
\end{cousinweb}

\begin{grammarlens}[title={What you've built}]
One more word for describing what you lived and felt.
\end{grammarlens}

\subsection*{Your turn}
\begin{itemize}
\raggedright
  \item \emph{Say it:} \emph{jān}
  \item \emph{Say it:} \emph{jān}, once more
  \item \emph{Say it:} \emph{sā̃je jān gāmmā̃ āvī.} — In the evening the groom's wedding party arrived in the village.
  \item \emph{Recall:} say the Gujarati for a wedding invitation card, then the Gujarati for the turmeric rite, then say \emph{jān} again
\end{itemize}

\begin{tcolorbox}[breakable,colback=teal!4,colframe=teal!35!black,title={Before you move on}]
What does \gu{જાન} mean? (\textquotedblleft{}The groom's wedding party\textquotedblright{}.) Say it once more.
\end{tcolorbox}
\section[\texorpdfstring{\gu{વરઘોડો} (varghoḍo)}{વરઘોડો (varghoḍo)}]{\gu{વરઘોડો} (varghoḍo) — the groom's procession}

\label{lesson:GU-C497-varghodo}



\begin{quote}
Before the new one: say the Gujarati for the turmeric rite, then the Gujarati for the groom's wedding party.
\end{quote}

\subsection*{You'll want to know: \gu{વરઘોડો}}
\textbf{\gu{વરઘોડો}} — \emph{varghoḍo} — \textquotedblleft{}the groom's procession\textquotedblright{}.

\begin{quote}
\gu{વરઘોડામાં} \gu{બધા} \gu{નાચ્યા}.
\end{quote}

\begin{quote}
\emph{varghoḍāmā̃ badhā nācyā.}
\end{quote}

\begin{quote}
Everyone danced in the groom's procession.
\end{quote}

A masculine noun.

\begin{cousinweb}
\textbf{\gu{વરઘોડો}} (\emph{varghoḍo}) is built from \emph{var}, \textquotedblleft{}bridegroom\textquotedblright{}, which CDIAL lists under Sanskrit \emph{vara}, \textquotedblleft{}suitor, bridegroom\textquotedblright{} (CDIAL 11309), and \textbf{\gu{ઘોડો}} (\emph{ghoḍo}), \textquotedblleft{}a horse\textquotedblright{} (CDIAL 4516): the groom rides to the wedding.
\end{cousinweb}

\begin{grammarlens}[title={What you've built}]
One more word for describing what you lived and felt.
\end{grammarlens}

\subsection*{Your turn}
\begin{itemize}
\raggedright
  \item \emph{Say it:} \emph{varghoḍo}
  \item \emph{Say it:} \emph{varghoḍo}, once more
  \item \emph{Say it:} \emph{varghoḍāmā̃ badhā nācyā.} — Everyone danced in the groom's procession.
  \item \emph{Recall:} say the Gujarati for the turmeric rite, then the Gujarati for the groom's wedding party, then say \emph{varghoḍo} again
\end{itemize}

\begin{tcolorbox}[breakable,colback=teal!4,colframe=teal!35!black,title={Before you move on}]
What does \gu{વરઘોડો} mean? (\textquotedblleft{}The groom's procession\textquotedblright{}.) Say it once more.
\end{tcolorbox}
\section[\texorpdfstring{\gu{પરણાવવું} (parṇāvvũ)}{પરણાવવું (parṇāvvũ)}]{\gu{પરણાવવું} (parṇāvvũ) — to marry off; to give in marriage}

\label{lesson:GU-C497-parnavvu}



\begin{quote}
Before the new one: say the Gujarati for the groom's wedding party, then the Gujarati for the groom's procession.
\end{quote}

\subsection*{You'll want to know: \gu{પરણાવવું}}
\textbf{\gu{પરણાવવું}} — \emph{parṇāvvũ} — \textquotedblleft{}to marry off; to give in marriage\textquotedblright{}.

\begin{quote}
\gu{તેઓ} \gu{આવતા} \gu{વરસે} \gu{દીકરીને} \gu{પરણાવશે}.
\end{quote}

\begin{quote}
\emph{teo āvtā varse dīkrīne parṇāvśe.}
\end{quote}

\begin{quote}
They will marry off their daughter next year.
\end{quote}

Past, with \gu{તેણે} (he or she …): \gu{તેણે} \gu{પોતાના} \gu{સંતાનને} \gu{પરણાવ્યું}, \emph{teṇe potānā santānne parṇāvyũ}.

\begin{cousinweb}
\textbf{\gu{પરણાવવું}} (\emph{parṇāvvũ}) is inherited from Sanskrit \emph{pariṇāpayati}, \textquotedblleft{}causes a man to marry\textquotedblright{}; CDIAL lists the Gujarati verb, \textquotedblleft{}to give in marriage\textquotedblright{} (CDIAL 7821).
\end{cousinweb}

\begin{grammarlens}[title={What you've built}]
One more word for describing what you lived and felt.
\end{grammarlens}

\subsection*{Your turn}
\begin{itemize}
\raggedright
  \item \emph{Say it:} \emph{parṇāvvũ}
  \item \emph{Say it:} \emph{parṇāvvũ}, once more
  \item \emph{Say it:} \emph{teo āvtā varse dīkrīne parṇāvśe.} — They will marry off their daughter next year.
  \item \emph{Recall:} say the Gujarati for the groom's wedding party, then the Gujarati for the groom's procession, then say \emph{parṇāvvũ} again
\end{itemize}

\begin{tcolorbox}[breakable,colback=teal!4,colframe=teal!35!black,title={Before you move on}]
What does \gu{પરણાવવું} mean? (\textquotedblleft{}To marry off; to give in marriage\textquotedblright{}.) Say it once more.
\end{tcolorbox}
